\documentclass[aspectratio=169]{beamer}

% Theme académique professionnel
\usetheme{Madrid}
\usecolortheme{default}

% Packages
\usepackage{amsmath,amsfonts}
\usepackage[french]{babel}

% Informations
\title{Time Series}
\subtitle{Decomposition: Trend Estimation}
\author{Mr. Aymen Ben Brik}

\institute{
Time Series Course 1GAMMA
}

\date{\today}


\usepackage{listings}
\usepackage{xcolor}

\lstset{
language=Python,
basicstyle=\ttfamily\small,
keywordstyle=\color{blue},
commentstyle=\color{gray},
stringstyle=\color{red},
showstringspaces=false,
frame=single,
breaklines=true
}



\begin{document}

% Slide titre
\begin{frame}
\titlepage
\end{frame}

% Objectifs
\begin{frame}{Session Objectives}

By the end of this session, you will be able to:

\begin{itemize}
\item Understand time series decomposition
\item Understand the trend component
\item Estimate and remove trend from time series
\item Apply moving average methods
\item Apply parametric regression methods
\item Understand non-parametric regression methods
\item Select appropriate trend estimation methods
\end{itemize}

\end{frame}

% Basics
\begin{frame}{Time Series: Basics}

Time series data:

\[
Y_t, \quad t = 1,2,\dots,T
\]

Model decomposition:

\[
Y_t = T_t + S_t + \varepsilon_t
\]

where:

\begin{itemize}

\item $T_t$ : trend component

\item $S_t$ : seasonal component

\item $\varepsilon_t$ : stationary component

\end{itemize}

Goal:

\begin{itemize}

\item Estimate $T_t$ and $S_t$

\item Remove them to obtain stationary process

\end{itemize}

\end{frame}

% Trend estimation overview
\begin{frame}{Trend Estimation}

Goal: eliminate trend component

Approaches:

\begin{itemize}

\item Estimate trend and subtract it

\item Difference the data

\end{itemize}

Methods:

\begin{itemize}

\item Moving Average

\item Parametric Regression

\item Non-Parametric Regression

\end{itemize}

\end{frame}

% Moving average
\begin{frame}{Trend Estimation: Moving Average}

Trend estimate:

\[
\hat{T}_t =
\frac{1}{2d +1}
\sum_{i=-k}^{k} Y_{t+i}
\]

where:

\begin{itemize}

\item $d$ : window size

\item used to smooth time series

\end{itemize}

Properties:

\begin{itemize}

\item Simple

\item Effective smoothing

\item Removes noise

\end{itemize}
\end{frame}

\begin{frame}[fragile]{Example: Moving Average in R}

Example: Estimate trend using moving average

\begin{lstlisting}[language=R]

# Generate example time series
t <- 1:100
y <- 0.5*t + rnorm(100,0,5)

# Moving average with window size 5
trend <- filter(y, rep(1/5,5), sides=2)

# Plot
plot(t,y,type="l",col="gray",
     main="Moving Average Trend Estimation")
lines(t,trend,col="blue",lwd=2)

\end{lstlisting}

\end{frame}
\begin{frame}[fragile]{Example: Moving Average in Python}

\begin{lstlisting}
import numpy as np
import matplotlib.pyplot as plt
import pandas as pd

# Generate time series
t = np.arange(100)
y = 0.5*t + np.random.normal(0,5,100)

# Moving average
trend = pd.Series(y).rolling(window=5).mean()

# Plot
plt.plot(t,y,color='gray')
plt.plot(t,trend,color='blue',linewidth=2)
plt.title("Moving Average Trend Estimation")
plt.show()
\end{lstlisting}

\end{frame}



% Parametric regression
\begin{frame}{Trend Estimation: Parametric Regression}

Assume polynomial model:

\[
T_t =
\beta_0
+
\beta_1 t
+
\beta_2 t^2
+
\dots
+
\beta_p t^p
\]

Estimation using regression:

\begin{itemize}

\item Linear trend ($p=1$)

\item Quadratic trend ($p=2$)

\end{itemize}

Model selection used to choose order.

\end{frame}

\begin{frame}[fragile]{Example: Linear Trend Estimation in R}

\begin{lstlisting}[language=R]

# Linear regression
model <- lm(y ~ t)

# Estimated trend
trend <- predict(model)

# Plot
plot(t,y,type="l",col="gray")
lines(t,trend,col="red",lwd=2)

\end{lstlisting}

\end{frame}

\begin{frame}[fragile]{Example: Linear Trend Estimation in Python}

\begin{lstlisting}

from sklearn.linear_model import LinearRegression

# reshape for sklearn
X = t.reshape(-1,1)

model = LinearRegression()
model.fit(X,y)

trend = model.predict(X)

plt.plot(t,y,color='gray')
plt.plot(t,trend,color='red',linewidth=2)
plt.title("Linear Trend Estimation")
plt.show()

\end{lstlisting}

\end{frame}



% Non-parametric regression
\begin{frame}{Trend Estimation: Non-Parametric Regression}

Kernel regression:

\[
\hat{T}_t =
\sum w_i Y_i
\]

Local polynomial regression:

\begin{itemize}

\item Fit polynomial locally

\item Flexible

\item Accurate

\end{itemize}

Other methods:

\begin{itemize}

\item Splines

\item Wavelets

\end{itemize}

\end{frame}

% Method comparison
\begin{frame}{Method Selection}

Preferred methods:

\begin{itemize}

\item Local polynomial regression

\item Splines regression

\end{itemize}

Advantages:

\begin{itemize}

\item Better accuracy

\item Handles boundaries well

\item Flexible modeling

\end{itemize}

Choice depends on smoothness of trend.

\end{frame}

\begin{frame}[fragile]{Example: Local Regression (LOESS) in R}

Local regression estimates the trend using local polynomial fitting.

\begin{lstlisting}[language=R]
# Generate example time series
t <- 1:100
y <- 0.5*t + 10*sin(t/10) + rnorm(100,0,5)
# Local regression using LOESS
model <- loess(y ~ t, span=0.2)
# Estimated trend
trend <- predict(model)
# Plot results
plot(t, y, type="l", col="gray",
     main="Local Regression (LOESS) Trend Estimation",
     xlab="Time", ylab="Value")
lines(t, trend, col="blue", lwd=2)
legend("topleft",
       legend=c("Original Series","LOESS Trend"),
       col=c("gray","blue"),lwd=2)

\end{lstlisting}

\end{frame}

\begin{frame}[fragile]{Example: Non-Parametric Trend (Python)}

\begin{lstlisting}

from statsmodels.nonparametric.smoothers_lowess import lowess

trend = lowess(y,t,frac=0.1)

plt.plot(t,y,color='gray')
plt.plot(trend[:,0],trend[:,1],
         color='green',linewidth=2)
plt.title("LOWESS Trend Estimation")
plt.show()

\end{lstlisting}

\end{frame}


% Summary
\begin{frame}{Summary}

In this session, we covered:

\begin{itemize}

\item Time series decomposition

\item Trend component

\item Moving average estimation

\item Parametric regression

\item Non-parametric regression

\end{itemize}

\end{frame}

\end{document}